\documentclass[11pt,a4paper]{article}
\usepackage[utf8]{inputenc}
\usepackage{amsmath,amssymb}
\usepackage{geometry}
\geometry{margin=2.5cm}
\usepackage{listings}
\usepackage{xcolor}
\usepackage{hyperref}

\lstset{
  language=C++,
  basicstyle=\ttfamily\small,
  keywordstyle=\color{blue}\bfseries,
  commentstyle=\color{green!60!black},
  stringstyle=\color{red!70!black},
  numbers=left,
  numberstyle=\tiny\color{gray},
  numbersep=8pt,
  frame=single,
  breaklines=true,
  tabsize=4,
  showstringspaces=false
}

\title{IOI 1998 -- Contact}
\author{Solution}
\date{}

\begin{document}
\maketitle

\section{Problem Statement}

Given a binary string $S$ of length $n$ and integers $a$, $b$, $k$
($1 \le a \le b \le 12$, $k \le 50$), find the most frequently occurring bit
patterns of length between $a$ and $b$ inclusive.

Output the top $k$ frequency groups (in decreasing order of frequency).
Within each group, patterns are sorted by length (shorter first), then by
numeric value (smaller first). Each pattern is printed in binary.

\medskip
\noindent\textbf{Constraints:} $|S| \le 200{,}000$.

\section{Solution Approach}

\subsection{Counting All Substrings}

Since pattern lengths are at most 12, each pattern fits in a 12-bit integer
(values 0 to $2^{12} - 1 = 4095$). We maintain a 2D frequency table
$\mathtt{cnt}[\ell][v]$ for length $\ell$ and value $v$.

For each starting position $i$, build the pattern incrementally:
\begin{verbatim}
  val = 0
  for len = 1 to min(b, n - i):
      val = val * 2 + S[i + len - 1]
      if len >= a:
          cnt[len][val]++
\end{verbatim}

\subsection{Sorting and Output}

Collect all (frequency, length, value) triples with positive frequency.
Sort by frequency descending, then length ascending, then value ascending.
Output the top $k$ frequency groups.

\section{C++ Solution}

\begin{lstlisting}
#include <cstdio>
#include <cstring>
#include <vector>
#include <algorithm>
using namespace std;

char buf[200005];
int cnt[13][4096]; // cnt[len][val]

int main() {
    int a, b, k;
    scanf("%d %d %d", &a, &b, &k);

    // Read binary string (may span multiple lines)
    int n = 0;
    char line[100];
    while (scanf("%s", line) == 1) {
        for (int i = 0; line[i]; i++)
            buf[n++] = line[i] - '0';
    }

    memset(cnt, 0, sizeof(cnt));

    // Count all patterns of length a..b
    for (int i = 0; i < n; i++) {
        int val = 0;
        for (int len = 1; len <= b && i + len <= n; len++) {
            val = val * 2 + buf[i + len - 1];
            if (len >= a)
                cnt[len][val]++;
        }
    }

    // Collect patterns with positive frequency
    struct Pattern {
        int freq, len, val;
    };
    vector<Pattern> patterns;
    for (int len = a; len <= b; len++)
        for (int v = 0; v < (1 << len); v++)
            if (cnt[len][v] > 0)
                patterns.push_back({cnt[len][v], len, v});

    // Sort: frequency descending, then length ascending, then value ascending
    sort(patterns.begin(), patterns.end(), [](const Pattern& x, const Pattern& y) {
        if (x.freq != y.freq) return x.freq > y.freq;
        if (x.len != y.len) return x.len < y.len;
        return x.val < y.val;
    });

    // Output top k frequency groups
    int groups = 0, i = 0;
    while (i < (int)patterns.size() && groups < k) {
        int curFreq = patterns[i].freq;
        printf("%d\n", curFreq);

        int lineCount = 0;
        while (i < (int)patterns.size() && patterns[i].freq == curFreq) {
            if (lineCount > 0) printf(" ");
            // Print binary representation with correct number of digits
            int len = patterns[i].len;
            int val = patterns[i].val;
            char bits[13];
            for (int j = len - 1; j >= 0; j--) {
                bits[j] = '0' + (val & 1);
                val >>= 1;
            }
            bits[len] = '\0';
            printf("%s", bits);
            lineCount++;
            i++;
        }
        printf("\n");
        groups++;
    }

    return 0;
}
\end{lstlisting}

\section{Complexity Analysis}

\begin{itemize}
  \item \textbf{Time complexity:} $O(n \cdot b)$ for counting, where $b \le 12$.
        Effectively $O(n)$. Sorting the patterns takes $O(P \log P)$ where
        $P \le (b - a + 1) \cdot 2^b \le 12 \cdot 4096 = 49{,}152$, which is negligible.
  \item \textbf{Space complexity:} $O(b \cdot 2^b + n)$: the frequency table plus the
        input string.
\end{itemize}

\end{document}
